\section{Introducción}

En el campo de la cardiologia debido a la complejidad y gravedad de las enfermades se hace de especial importancia el uso de herramientas capaces de leer y analizar datos clinicos de los pacientes de manera que se pueda advertir de posibles riesgos para su salud, sabiendo como esta tarea puede encontrar solucion en la programacion y las estructuras de datos como grupo vemos de especial interes hacer uso de una implementacion de lo visto en el curso
para clasificar los datos clinicos de un paciente de forma que se puedan leer y analizar de manera eficiente por parte del personal de salud.

Para este fin en nuestro proyecto se estudiara la lectura, analisis y graficacion de señales ECG (obtenidas a base de un electrocardiograma) para ser capaces de entregar un reporte en el que se puedan identificar patrones anomalos.

\subsection{Definicion del proyecto}
La idea central del proyecto consiste en por medio de las implementaciones estudiadas en el curso de estructuras de datos manejar los datos del ECG y detectar anomalias, esto puede ser de gran ayuda pues las estructuras de datos vistas pueden hacer mas eficiente el costo de operacion y por tanto facilitar el trabajo para el personal medico. 
\textbf{[Describir con precisión la situación que se desea resolver, sus
causas, restricciones y consecuencias.]}

\subsection{Objetivos}

\begin{itemize}
    \item \textbf{Ser capaz de leer registros de la base de datos MIT-BIH Arrhythmia Database incluyendo los archivos de registro y las anotaciones del cardiologo.}
    \item \textbf{Usar estructuras de datos capaces de permitir un recibiendo y envio eficiente de datos para el procesamiento y filtrado de informacion}
    \item \textbf{Entregar un resultado entendible que permita analizar los datos}
\end{itemize}